% =====================================================================
\chapter{Demostración directa y por contrarrecíproco}
\label{tema:3}
% =====================================================================

La mayoría de los teoremas tienen la forma $\forall x\,(P(x)\Rightarrow Q(x))$. Este tema presenta las dos técnicas básicas para demostrarlos: la demostración directa, que parte de $P(x)$ y llega a $Q(x)$, y la demostración por contrarrecíproco, que parte de $\neg Q(x)$ y llega a $\neg P(x)$. Como banco de pruebas se usa la aritmética de los enteros (paridad, divisibilidad y congruencias), igual que en los capítulos~4 y~5 de~\cite{hammack}, la §3.2 de~\cite{velleman}, las §1.5–1.6 de~\cite{mcs} y el primer \emph{problem set} de CS103~\cite{cs103}.

\section{Definiciones de partida}

Una demostración solo puede usar lo que está definido o demostrado antes. Se admiten como conocidas las propiedades algebraicas habituales de $\Z$ (la suma, la resta y el producto de enteros son enteros; asociatividad, conmutatividad, distributividad) y de $\R$ (incluidas las reglas de las desigualdades), y se fijan las definiciones siguientes~\cite[§4.2]{hammack}.

\begin{definicion}[Paridad]\label{def:paridad}
Un entero $n$ es \emph{par} si existe $k\in\Z$ con $n=2k$, e \emph{impar} si existe $k\in\Z$ con $n=2k+1$. Dos enteros tienen la \emph{misma paridad} si ambos son pares o ambos impares.
\end{definicion}

\begin{definicion}[Divisibilidad]\label{def:divide}
Sean $a,b\in\Z$. Se dice que $a$ \emph{divide} a $b$, y se escribe $a\divides b$, si existe $c\in\Z$ con $b=ac$. Se dice también que $b$ es \emph{múltiplo} de $a$ y que $a$ es \emph{divisor} de $b$. Un entero $p\ge 2$ es \emph{primo} si sus únicos divisores positivos son $1$ y $p$; un entero $n\ge 2$ que no es primo es \emph{compuesto}.
\end{definicion}

\begin{teorema}[División euclídea]\label{teo:division}
Sean $a\in\Z$ y $b\in\N$. Existen enteros únicos $q$ y $r$ con $a=qb+r$ y $0\le r<b$.
\end{teorema}

La existencia se demuestra en el \cref{tema:5} (\cref{prop:division-existencia}) y la unicidad en el \cref{tema:6} (\cref{prop:division-unicidad}). Con $b=2$, el teorema dice que todo entero es par o impar, y no ambas cosas; este hecho se usará libremente.

\section{Demostración directa}

\begin{definicion}[Demostración directa]\label{def:directa}
Para demostrar $\forall x\,(P(x)\Rightarrow Q(x))$ de forma \emph{directa}, se toma un $x$ arbitrario, se supone $P(x)$ y, mediante definiciones, resultados previos y reglas de inferencia, se deduce $Q(x)$.
\end{definicion}

El esquema es siempre el mismo: «Sea $x$ [del universo] tal que $P(x)$. [\dots] Por tanto, $Q(x)$.» Antes de empezar conviene escribir qué significa la tesis según las definiciones, para saber a qué expresión hay que llegar~\cite[§4.3]{hammack}.

\begin{proposicion}\label{prop:impar-cuadrado}
Si $n$ es un entero impar, entonces $n^{2}$ es impar.
\end{proposicion}

\begin{proof}
Sea $n$ un entero impar. Por definición existe $k\in\Z$ con $n=2k+1$. Entonces
\[
n^{2}=4k^{2}+4k+1=2(2k^{2}+2k)+1.
\]
Como $m=2k^{2}+2k$ es un entero, $n^{2}=2m+1$ es impar.
\end{proof}

Obsérvese que la demostración introduce un nombre nuevo, $m$, para el entero que da la forma $2m+1$: el objetivo era encontrar \emph{algún} entero con esa propiedad, y la definición de impar no exige que sea el mismo $k$ de la hipótesis.

\begin{proposicion}[Propiedades de la divisibilidad]\label{prop:divisibilidad}
Sean $a,b,c,x,y\in\Z$.
\begin{enumerate}
  \item Si $a\divides b$ y $b\divides c$, entonces $a\divides c$.
  \item Si $a\divides b$ y $a\divides c$, entonces $a\divides(bx+cy)$.
\end{enumerate}
\end{proposicion}

\begin{proof}
(i) Supóngase $a\divides b$ y $b\divides c$. Existen $d,e\in\Z$ con $b=ad$ y $c=be$. Sustituyendo, $c=(ad)e=a(de)$ con $de\in\Z$, luego $a\divides c$.

(ii) Supóngase $a\divides b$ y $a\divides c$, y sean $d,e\in\Z$ con $b=ad$ y $c=ae$. Entonces $bx+cy=adx+aey=a(dx+ey)$, y $dx+ey\in\Z$, luego $a\divides(bx+cy)$.
\end{proof}

\begin{ejemplo}[Desigualdad con números reales]\label{ej:amgm}
Se demuestra que, si $x,y\in\R$ son positivos, entonces $\sqrt{xy}\le\frac{x+y}{2}$. Una forma de \emph{descubrir} la prueba es trabajar hacia atrás: la desigualdad equivale a $2\sqrt{xy}\le x+y$, es decir, a $0\le x-2\sqrt{xy}+y=(\sqrt x-\sqrt y)^{2}$, que es cierta. La demostración se \emph{redacta}, sin embargo, hacia delante: «Sean $x,y>0$. Como el cuadrado de un número real es no negativo, $0\le(\sqrt x-\sqrt y)^{2}=x-2\sqrt{xy}+y$. Sumando $2\sqrt{xy}$ y dividiendo por~2 se obtiene $\sqrt{xy}\le(x+y)/2$.» Cada paso de la redacción es una consecuencia de los anteriores, y no al revés.
\end{ejemplo}

\section{Demostración por contrarrecíproco}

\begin{definicion}[Demostración por contrarrecíproco]\label{def:contrarreciproco}
Para demostrar $\forall x\,(P(x)\Rightarrow Q(x))$ por \emph{contrarrecíproco}, se demuestra de forma directa $\forall x\,(\neg Q(x)\Rightarrow\neg P(x))$. Es legítimo porque ambos enunciados son lógicamente equivalentes (\cref{prop:leyes}(v)).
\end{definicion}

La técnica es útil cuando la negación de la tesis aporta información más manejable que la hipótesis. Un indicio típico es que la tesis afirme algo sobre un número ($n$ es par) y la hipótesis, algo sobre una expresión de ese número ($n^{2}$ es par)~\cite[§5.1]{hammack}.

\begin{proposicion}\label{prop:cuadrado-par}
Sea $n\in\Z$. Si $n^{2}$ es par, entonces $n$ es par.
\end{proposicion}

\begin{proof}
Se demuestra el contrarrecíproco: si $n$ no es par, entonces $n^{2}$ no es par. Si $n$ no es par, es impar, y por la \cref{prop:impar-cuadrado} $n^{2}$ es impar, luego no es par.
\end{proof}

Un intento directo habría partido de $n^{2}=2k$ y necesitaría extraer una raíz cuadrada, que no conduce a ninguna parte. El contrarrecíproco convierte el problema en el cálculo ya hecho.

\begin{ejemplo}\label{ej:contra-real}
Sea $x\in\R$. Si $x^{3}-x>0$, entonces $x>-1$. El contrarrecíproco dice: si $x\le -1$, entonces $x^{3}-x\le 0$. En efecto, $x^{3}-x=x(x-1)(x+1)$; para $x\le -1$ se tiene $x<0$, $x-1<0$ y $x+1\le 0$, y el producto de dos factores negativos y uno no positivo es no positivo.
\end{ejemplo}

\section{Congruencias}

\begin{definicion}[Congruencia]\label{def:congruencia}
Sea $n\in\N$. Dos enteros $a$ y $b$ son \emph{congruentes módulo $n$}, y se escribe $a\equiv b\pmod n$, si $n\divides(a-b)$.
\end{definicion}

Por la división euclídea, todo entero es congruente módulo $n$ con exactamente uno de $0,1,\dots,n-1$, su resto. Por ejemplo, $n$ es par si y solo si $n\equiv 0\pmod 2$.

\begin{proposicion}\label{prop:congruencias}
Sea $n\in\N$ y sean $a,b,c,d\in\Z$ con $a\equiv b\pmod n$ y $c\equiv d\pmod n$. Entonces
\begin{enumerate}
  \item $a+c\equiv b+d\pmod n$;
  \item $ac\equiv bd\pmod n$;
  \item $a^{k}\equiv b^{k}\pmod n$ para todo $k\in\N$.
\end{enumerate}
\end{proposicion}

\begin{proof}
Por hipótesis existen $s,t\in\Z$ con $a-b=ns$ y $c-d=nt$.
(i) $(a+c)-(b+d)=ns+nt=n(s+t)$.
(ii) $ac-bd=ac-bc+bc-bd=(a-b)c+b(c-d)=n(sc+bt)$.
(iii) Se aplica (ii) con $c=a$ y $d=b$ repetidamente: $a^{2}\equiv b^{2}$, $a^{3}=a^{2}a\equiv b^{2}b=b^{3}$, etc.; la forma rigurosa de este «etc.» es la inducción del \cref{tema:5}.
\end{proof}

\begin{ejemplo}\label{ej:resto-potencia}
Se calcula el resto de dividir $2^{100}$ entre~7. Como $2^{3}=8\equiv 1\pmod 7$, la \cref{prop:congruencias}(iii) da $2^{99}=(2^{3})^{33}\equiv 1^{33}=1$, y por (ii), $2^{100}=2^{99}\cdot 2\equiv 2\pmod 7$. El resto es~2.
\end{ejemplo}

\section{Bicondicionales y cadenas de equivalencias}

\begin{definicion}[Demostración de un bicondicional]
Para demostrar $P\Leftrightarrow Q$ se demuestran por separado $P\Rightarrow Q$ y $Q\Rightarrow P$, cada una con la técnica más conveniente (\cref{prop:leyes}(vii)). Para demostrar que varias afirmaciones $P_1,\dots,P_k$ son equivalentes basta una cadena cíclica $P_1\Rightarrow P_2\Rightarrow\dots\Rightarrow P_k\Rightarrow P_1$~\cite[§7.2]{hammack}.
\end{definicion}

\begin{proposicion}\label{prop:par-equiv}
Para $n\in\Z$, son equivalentes: (a) $n$ es par; (b) $n^{2}$ es par; (c) $n^{2}+2n$ es múltiplo de~4.
\end{proposicion}

\begin{proof}
(a)$\Rightarrow$(c): si $n=2k$, entonces $n^{2}+2n=4k^{2}+4k=4(k^{2}+k)$.
(c)$\Rightarrow$(b): si $n^{2}+2n=4m$, entonces $n^{2}=4m-2n=2(2m-n)$ es par.
(b)$\Rightarrow$(a): es la \cref{prop:cuadrado-par}.
\end{proof}

\begin{observacion}[Errores frecuentes]\label{obs:3-errores}
(a) Suponer la tesis: una «demostración» que empieza escribiendo lo que se quiere probar y llega a algo cierto no demuestra nada, salvo que cada paso sea reversible y se redacte en el sentido correcto (\cref{ej:amgm}).
(b) Reutilizar una letra: si $a$ y $b$ son impares, no puede escribirse $a=2k+1$ y $b=2k+1$ con el mismo $k$, pues eso supone $a=b$.
(c) Comprobar unos cuantos casos ($n=1,3,5$) no demuestra un enunciado universal.
(d) Demostrar el recíproco $Q\Rightarrow P$ en lugar del contrarrecíproco $\neg Q\Rightarrow\neg P$.
\end{observacion}

\begin{observacion}[Conexión con el curso]
Una función $f$ es inyectiva si $f(x)=f(y)\Rightarrow x=y$, lo que por contrarrecíproco equivale a $x\neq y\Rightarrow f(x)\neq f(y)$. Las dos formas se usan en el curso: en el módulo~1 se prueba de forma directa que un tokenizador $\tau$ con decodificador $\delta$ cumple $\tau(s)=\tau(s')\Rightarrow s=s'$, aplicando $\delta$ a ambos lados; en el módulo~12 se demuestra la inyectividad de la codificación UTF-8 de manera análoga. Las congruencias reaparecen en la aritmética binaria y hexadecimal del módulo~12 y en el análisis de la coma flotante del módulo~11.
\end{observacion}

\begin{lecturas}
\item Hammack~\cite{hammack}, cap.~4 «Direct Proof» (§4.1–4.3) y cap.~5 «Contrapositive Proof» (§5.1–5.2).
\item Velleman~\cite{velleman}, §3.1 «Proof Strategies» y §3.2 «Proofs Involving Negations and Conditionals».
\item Lehman, Leighton y Meyer~\cite{mcs}, §1.5 «Proving an Implication» y §1.6 «Proving an “If and Only If”».
\item Stanford CS103~\cite{cs103}, clase~1 «Mathematical Proofs» y la guía «Guide to Proofs».
\end{lecturas}

\section*{Ejercicios del tema 3}
\addcontentsline{toc}{section}{Ejercicios del tema 3}

\begin{ejercicio}{1}{eje:3-hipotesis}
Para cada enunciado, identifique hipótesis y tesis, escriba el contrarrecíproco y diga qué técnica parece más sencilla: (a) «si $a\divides b$ y $b\divides c$, entonces $a\divides c$»; (b) «si $n^{2}$ es impar, entonces $n$ es impar»; (c) «si $x\in\R$ y $x^{2}+x<0$, entonces $x<0$»; (d) «si $xy$ es irracional, entonces $x$ o $y$ es irracional».
\end{ejercicio}

\begin{ejercicio}{1}{eje:3-esquema}
Explique qué está mal en esta demostración de «si $n$ es par, entonces $n^{2}$ es múltiplo de~4»: «Sea $n^{2}=4k$. Entonces $n=2\sqrt{k}$, que es par. Por tanto, el resultado es cierto.» Redáctela correctamente.
\end{ejercicio}

\begin{ejercicio}{3}{eje:3-cubo}
Demuestre que, si $x$ es un entero impar, entonces $x^{3}$ es impar.
\fuente{adaptado de Hammack~\cite{hammack}, cap.~4, ejercicio~2.}
\end{ejercicio}

\begin{ejercicio}{3}{eje:3-polinomio}
Demuestre que, si $a$ es un entero impar, entonces $a^{2}+3a+5$ es impar.
\fuente{adaptado de Hammack~\cite{hammack}, cap.~4, ejercicio~3.}
\end{ejercicio}

\begin{ejercicio}{3}{eje:3-cuadrados}
Sean $a,b\in\Z$. Demuestre que, si $a\divides b$, entonces $a^{2}\divides b^{2}$.
\fuente{adaptado de Hammack~\cite{hammack}, cap.~4, ejercicio~7.}
\end{ejercicio}

\begin{ejercicio}{3}{eje:3-producto}
Sean $a,b,c,d\in\Z$. Demuestre que, si $a\divides b$ y $c\divides d$, entonces $ac\divides bd$.
\fuente{adaptado de Hammack~\cite{hammack}, cap.~4, ejercicio~11.}
\end{ejercicio}

\begin{ejercicio}{3}{eje:3-desigualdad}
Demuestre que, si $x\in\R$ y $0<x<4$, entonces $\dfrac{4}{x(4-x)}\ge 1$.
\fuente{adaptado de Hammack~\cite{hammack}, cap.~4, ejercicio~12.}
\end{ejercicio}

\begin{ejercicio}{3}{eje:3-diferencia}
Demuestre que todo entero impar es diferencia de los cuadrados de dos enteros.
\fuente{adaptado de Hammack~\cite{hammack}, cap.~4, ejercicio~26.}
\end{ejercicio}

\begin{ejercicio}{3}{eje:3-contra-cubo}
Sea $n\in\Z$. Demuestre por contrarrecíproco que, si $n^{3}+5$ es impar, entonces $n$ es par.
\end{ejercicio}

\begin{ejercicio}{3}{eje:3-pitagoras}
Sean $a,b,c\in\Z$. Demuestre que, si $a^{2}+b^{2}=c^{2}$, entonces al menos uno de los tres números es par.
\fuente{adaptado de Stanford CS103~\cite{cs103}, \emph{Problem Set}~1, problema~5 «Writing Proofs by Contrapositive».}
\end{ejercicio}

\begin{ejercicio}{3}{eje:3-equivalencia}
Sea $k\in\N$ y defínase, para $x,y\in\Z$, $x\equiv_k y$ si existe $q\in\Z$ con $x=y+kq$. Demuestre que la relación es (a) reflexiva: $x\equiv_k x$; (b) simétrica: $x\equiv_k y\Rightarrow y\equiv_k x$; (c) transitiva: $x\equiv_k y$ y $y\equiv_k z$ implican $x\equiv_k z$.
\fuente{adaptado de Stanford CS103~\cite{cs103}, \emph{Problem Set}~1, problema~8.}
\end{ejercicio}

\begin{ejercicio}{3}{eje:3-real}
Sean $x,y\in\R$. Demuestre que, si $y^{3}+yx^{2}\le x^{3}+xy^{2}$, entonces $y\le x$.
\end{ejercicio}

\begin{ejercicio}{3}{eje:3-bicond}
Sea $n\in\Z$. Demuestre que $4\divides n^{2}$ si y solo si $n$ es par.
\end{ejercicio}

\begin{ejercicio}{2}{eje:3-restos}
Usando la \cref{prop:congruencias}, calcule sin calculadora: (a) el resto de dividir $3^{100}$ entre~7; (b) la última cifra de $7^{2026}$; (c) el resto de dividir $5^{40}+6^{40}$ entre~11.
\end{ejercicio}

\begin{ejercicio}{2}{eje:3-division}
(a) Calcule el cociente y el resto de la división euclídea de $2026$ entre~$7$ y de $-17$ entre~$5$. (b) Escriba la definición que hay que comprobar para decidir si $7\divides 2023$ y $7\divides 2026$, y decídalo. (c) Calcule el resto de $2026$ módulo~$9$ y compruebe que coincide con el resto de la suma de sus cifras.
\end{ejercicio}
